% chapters/05/section-04.tex
\section{Menyinkronkan Posisi Pecahan dengan \code{Rect}}
\label{sec:selaras-rect}

Sampai di sini posisi objek berupa vektor dengan nilai pecahan, sedangkan
seluruh perhitungan gerak berjalan dengan lancar. Persoalannya muncul ketika
posisi itu hendak dipakai untuk menggambar.

\subsection{Dua Tipe Data yang Berbeda}

\code{pygame.Vector2} menyimpan nilai bertipe pecahan, misalnya 312,47.
\code{pygame.Rect} menyimpan nilai bertipe bilangan bulat, dan memang
dirancang demikian karena letak piksel pada layar tidak memiliki pecahan:
piksel ke-312 tepat berada di tempatnya, dan tidak ada piksel ke-312,47.

Karena itu, posisi yang berupa pecahan harus diselaraskan terlebih dahulu
sebelum dipakai menggambar. Penyelarasan itu dilakukan setiap bingkai, tepat
sebelum penggambaran.

\begin{kodekeluaran}
player_rect.center = (round(position.x), round(position.y))
pygame.draw.rect(screen, PLAYER_COLOR, player_rect)
\end{kodekeluaran}

Pembulatan dilakukan karena posisi sesungguhnya tetap disimpan pada
\code{position} sebagai pecahan. \code{player_rect} hanyalah salinan yang
dibulatkan untuk keperluan menggambar dan---pada Bab~VII nanti---untuk
memeriksa tabrakan.

\subsection{Posisi Tetap Disimpan sebagai Pecahan}

Aturan yang perlu Anda pegang adalah sebagai berikut.

\begin{itemize}
  \item \textbf{Posisi yang sebenarnya disimpan sebagai pecahan}
        (\code{Vector2}), karena hanya dengan begitu gerak dapat berlangsung
        halus. Pada kecepatan 20 piksel per detik dan laju 60 bingkai per
        detik, perpindahan setiap bingkainya hanya 0,33 piksel; nilai itu harus
        tersimpan agar dapat terakumulasi.
  \item \textbf{\code{Rect} dipakai untuk menggambar dan memeriksa tabrakan},
        dan diselaraskan dari posisi pecahan setiap bingkai. \code{Rect} bukan
        tempat menyimpan posisi.
\end{itemize}

Akibat apabila aturan itu dilanggar diperlihatkan pada
Bagian~\ref{sec:contoh-bab-5}. Objek yang posisinya disimpan sebagai bilangan
bulat tidak akan bergerak sama sekali pada kecepatan rendah. Perpindahan 0,33
piksel terpotong menjadi nol pada setiap bingkai, dan nilai yang terpotong itu
hilang.

\subsection{Membulatkan atau Memotong}

Pilihan antara \code{round()} dan \code{int()} pada penyelarasan di atas
memengaruhi letak objek yang terlihat. Keduanya memberi hasil yang sama untuk
nilai pecahan di bawah 0,5, tetapi berbeda untuk nilai di atasnya.

\begin{kodekeluaran}
int(312.9)      hasilnya 312      (selalu ke arah nol)
round(312.9)    hasilnya 313      (ke bilangan bulat terdekat)
\end{kodekeluaran}

Karena posisi yang benar adalah 312,9, maka nilai yang lebih dekat adalah 313.
Memakai \code{int()} akan membuat objek tampak tertinggal sejauh hampir satu
piksel secara sistematis. Selisih satu piksel itu tidak terlihat pada objek
tunggal, tetapi menjadi terlihat sebagai gerak yang tersendat ketika beberapa
objek bergerak bersama-sama, misalnya pada Bab~VII nanti.

\begin{catatabox}
Perbedaan antara membulatkan dan memotong juga menjelaskan mengapa posisi
tidak boleh disimpan di dalam \code{Rect}. Apabila posisi disimpan sebagai
bilangan bulat, membulatkan atau memotong tidak lagi menjadi pilihan: nilai
pecahannya sudah hilang sebelum sempat terakumulasi.
\end{catatabox}

\subsection{Mengapa Pemisahan Ini Bermanfaat}

Memisahkan ``posisi yang sebenarnya'' dari ``posisi yang digambar'' membawa
tiga manfaat yang akan terasa pada bab-bab berikutnya.

\begin{enumerate}
  \item \textbf{Gerak tetap halus pada kecepatan rendah.} Perpindahan pecahan
        tetap terakumulasi walaupun hasil penggambarannya baru berubah
        sesekali.
  \item \textbf{Perhitungan tabrakan tetap dapat dilakukan.} Pemeriksaan
        tabrakan memerlukan \code{Rect} bertipe bilangan bulat, dan
        \code{Rect} itu tersedia setiap bingkai tanpa mengorbankan ketelitian
        posisi.
  \item \textbf{Perubahan laju bingkai tidak mengubah posisi akhir.} Karena
        posisi yang sebenarnya tidak bergantung pada pembulatan, dua komputer
        dengan laju bingkai berbeda akan menghasilkan posisi akhir yang sama,
        walaupun gambar yang ditampilkan sesaat berbeda.
\end{enumerate}
